\chapter{Antisocial Personality Disorder}
\index{Antisocial Personality Disorder}

\section*{Definition and Overview}
Antisocial personality disorder (ASPD) is a mental health condition characterised by a persistent and pervasive pattern of disregard for, and violation of, the rights of others. People with ASPD repeatedly break rules and laws, deceive and manipulate others, act impulsively, and show little or no remorse for the harm they cause. The pattern begins in childhood or adolescence -- most often as conduct disorder -- and continues into adult life.

ASPD belongs to the group of personality disorders sometimes described as the ``dramatic'' or cluster B conditions, alongside borderline and narcissistic personality disorder. It is important to recognise that ASPD is not the same as criminality: not all people with the condition come into contact with the law, and not all offenders have ASPD. Nevertheless, criminal behaviour, substance misuse, and unstable relationships are common, and the condition is associated with substantial harm to others and to the person themselves.

\section*{Epidemiology}
ASPD is estimated to affect roughly 1 to 3 per cent of adults, making it one of the more common personality disorders. It is diagnosed more often in men than in women, at a ratio of about three to one, and it is over-represented in forensic populations: prevalence among prisoners is several times higher than in the general community.

The condition typically becomes recognisable in adolescence, with conduct disorder in childhood a frequent forerunner. It is more common among people with substance use disorders, unstable family backgrounds, and lower socioeconomic circumstances. Because people with ASPD rarely present to health services voluntarily, the true prevalence is likely undercounted, and many are first identified through the justice system, emergency services, or family members.

\section*{Aetiology and Pathophysiology}
ASPD arises from a combination of genetic, biological, and environmental influences rather than a single cause:

\begin{itemize}
    \item \textbf{Genetics:} the disorder runs in families, with heritability estimated at around 50 per cent
    \item \textbf{Brain function:} imaging studies show reduced activity in the prefrontal cortex, the region involved in impulse control, planning, and social judgement, together with altered emotional responses in the amygdala
    \item \textbf{Childhood environment:} physical or sexual abuse, neglect, unstable parenting, and exposure to domestic violence increase risk substantially
    \item \textbf{Parenting style:} inconsistent or harsh discipline, and parental criminality or substance misuse, are associated with the development of the disorder
    \item \textbf{Temperament:} some children are born with traits such as fearlessness, sensation-seeking, and low sensitivity to punishment that make them more vulnerable
    \item \textbf{Neurobiology of reward and threat:} altered processing of threat and reward may explain the coldness, risk-taking, and difficulty learning from punishment
\end{itemize}

\section*{Clinical Features}
\begin{itemize}
    \item \textbf{Repeated law-breaking or rule-breaking} that begins in adolescence and continues into adulthood
    \item \textbf{Deceitfulness:} lying, using false identities, or conning others for personal gain
    \item \textbf{Impulsivity} and a persistent failure to plan ahead
    \item \textbf{Irritability and aggressiveness:} repeated physical fights or assaults
    \item \textbf{Reckless disregard for safety} of self or others, including dangerous driving and risk-taking
    \item \textbf{Consistent irresponsibility:} failing to sustain work, honour financial obligations, or care for dependants
    \item \textbf{Lack of remorse:} indifference to, or rationalisation of, hurt or harm caused to others
    \item \textbf{Superficial charm and manipulation:} exploiting relationships for personal advantage without genuine attachment
\end{itemize}

For a diagnosis in adulthood, there must be evidence of conduct disorder before the age of 15.

\section*{Differential Diagnosis}
\begin{itemize}
    \item \textbf{Borderline personality disorder:} emotional instability, intense and unstable relationships, and self-harm, with more capacity for remorse
    \item \textbf{Narcissistic personality disorder:} grandiosity and entitlement rather than rule-breaking and impulsivity
    \item \textbf{Conduct disorder:} the equivalent pattern in adolescents; ASPD is only diagnosed from age 18
    \item \textbf{Bipolar disorder:} impulsive and aggressive behaviour during manic episodes, which is episodic rather than persistent
    \item \textbf{Substance use disorders:} intoxication or dependence can closely mimic antisocial behaviour
    \item \textbf{Attention deficit hyperactivity disorder:} impulsivity and rule-breaking without the characteristic disregard for others
    \item \textbf{Acquired frontal lobe damage:} head injury, stroke, or tumours affecting the frontal lobes can change personality
    \item \textbf{Antisocial behaviour from social deprivation:} should be distinguished from a persistent personality pattern
\end{itemize}

\section*{When to Seek Medical Care}
See a GP or mental health professional if you are concerned about your own behaviour, or that of a family member, particularly if it is causing harm, conflict, or legal problems. Early engagement improves outcomes, and co-existing conditions such as depression, anxiety, and substance misuse should be treated at the same time, since they drive much of the risk.

\textbf{Contact your local emergency services for:}
\begin{itemize}
    \item Threats of serious violence toward another person
    \item Threats of suicide or self-harm
    \item Any situation in which you or someone else is in immediate physical danger
    \item Severe agitation, confusion, or an acute mental health crisis
\end{itemize}

\section*{Diagnosis and Investigations}
\begin{itemize}
    \item \textbf{Clinical interview:} a detailed psychiatric history covering behaviour across work, relationships, finances, and the law, and the presence of conduct disorder in childhood
    \item \textbf{Collateral information:} accounts from family, partners, employers, or legal and medical records, which are often essential because self-report is unreliable
    \item \textbf{Mental state examination:} assessment of mood, thought, perception, and risk to self and others
    \item \textbf{Standardised tools:} personality and behavioural rating scales used by experienced clinicians to support diagnosis
    \item \textbf{Physical examination and blood tests:} to exclude medical, neurological, or substance-related causes of the behaviour
    \item \textbf{Assessment of co-morbidity:} structured screening for depression, anxiety, post-traumatic stress, and substance use disorders
    \item \textbf{Risk assessment:} structured judgement of risk of violence, self-harm, and exploitation
\end{itemize}

\section*{Management}
\begin{itemize}
    \item \textbf{Psychological therapy:} cognitive behaviour therapy, mentalisation-based approaches, and anger or aggression management may help, although engagement and motivation are often limited
    \item \textbf{Treatment of co-morbidity:} addressing substance misuse and depression is among the most effective ways to reduce destructive behaviour
    \item \textbf{Medication:} there is no specific medicine for ASPD itself, but drugs may be used for co-existing conditions and for severe aggression, always under specialist supervision
    \item \textbf{Structured support:} forensic psychology and community teams with experience in personality disorder provide consistent, firm, and non-punitive care
    \item \textbf{Behavioural programmes:} problem-solving, social skills, and vocational programmes in custodial and community settings
    \item \textbf{Family and carer support:} education, safety planning, and support for partners and relatives affected by the behaviour
\end{itemize}

\section*{Complications}
\begin{itemize}
    \item Imprisonment and long-standing legal consequences
    \item Violence and harm to others, including family members
    \item Substance misuse and high-risk behaviour
    \item Suicide, accidents, and premature death
    \item Relationship breakdown, social isolation, and exploitation of others
    \item Financial instability, debt, and homelessness
    \item Poor engagement with healthcare, with increased rates of physical illness and reduced life expectancy
\end{itemize}

\section*{Prognosis and Outcomes}
ASPD is a lifelong personality pattern, but it does not remain constant. Impulsive, aggressive, and criminal behaviour tends to diminish with age, often becoming less prominent from around the fourth and fifth decades of life -- a phenomenon sometimes described as ``burnout''. Outcomes vary widely: some people stabilise with structure, treatment of co-morbidities, and social support, while others continue to cause significant harm.

Personality traits themselves rarely change, but harmful behaviours can improve substantially, particularly when substance misuse is treated and when consistent support is in place. Prognosis is better for people with higher intelligence, stronger social supports, and less severe childhood adversity; it is worse when alcohol or drug dependence dominates and when the person remains disengaged from services.

\section*{Prevention and Self-Care}
\begin{itemize}
    \item Support early intervention for children with conduct problems, including evidence-based parenting programmes
    \item Provide stable, consistent, and non-violent caregiving in childhood
    \item Treat substance misuse early in adolescents and adults
    \item Address childhood trauma and abuse through counselling and support
    \item Build structured routines, meaningful employment, and stable social connections
    \item Learn to identify and avoid situations that trigger impulsive or aggressive behaviour
    \item Seek help for co-existing depression, anxiety, or anger problems
    \item For family members: seek support and safety planning from mental health services and advocacy groups
\end{itemize}

\section*{Sources and Further Reading}
The following organisations provide reliable, up-to-date information on antisocial personality disorder for patients, students, and health professionals.

\begin{itemize}
    \item NHS. \textit{Information on antisocial personality disorder and its treatment.} https://www.nhs.uk/conditions/
    \item NICE. \textit{Guidelines on antisocial personality disorder and management.} https://www.nice.org.uk/
    \item Mayo Clinic. \textit{Overview of antisocial personality disorder, symptoms, and treatment.} https://www.mayoclinic.org/
    \item MedlinePlus. \textit{Health information on personality disorders.} https://medlineplus.gov/
    \item MSD Manual. \textit{Professional overview of antisocial personality disorder.} https://www.msdmanuals.com/
    \item World Health Organization. \textit{Resources on mental health and personality disorders.} https://www.who.int/
\end{itemize}
